\subsection{Overall System Architecture}
The system proposed in the paper is designed as a modular IoT-based Environmental
monitoring platform, consisting of an embedded sensing processor subsystem and a higher-level
telemetry processor. The embedded subsystem is responsible for sensor data acquisition, fan control
and monitoring, aswell as low-level hardware interaction.\\

The sensing subsystem is based on a PIC16F18124[INDSÆT REF] microcontroller, with its own
auxillary supply, connected to an ENS160 air-quality sensor and multiple PWM-controlled fans. 
The microcontroller continuously acquires environmental measurements and fan telemetry while
simultaneously regulating the fan speed. System information is transmitted to the Raspberry PI
through the UART communication interface.\\

The Raspberry Pi acts a edge gateway between the embedded hardware and the IoT infrastructure.
Incoming UART telemetry is processed by a python application and published to MQTT topics.
This telemetry is subsequently collected by telegraf to be stored in an InfluxDB time-series 
databased, later to be visualized by grafana. It is the handling done by telegraf that 
timestamps the data, as the microcontroller system, effectively can't determine absolute 
time, only that which have passed since boot. 

\subsection{PIC Firmware architecture}
